\scanpage{24}

\mspage{11}

Comme \unsure{tt|tout} autom d'une structure $2$-polyèdrale \inserted{(au sens strict)} connexe est connu quand on \unsure{connait|connaît} son effet sur un repère, on trouve le

\noindent\vrule\hspace{.5em}\parbox{.92\linewidth}{\uline{Théorème} Si $\Lambda=\mathbb{Z}$ ou $\mathbb{Z}/n\mathbb{Z}$ $(n\geq3)$, alors les hom verticaux dans \unsure{50)|(50)|(50),} sont des isomorphismes, i.e.}

\[
\left\{\begin{aligned}
\mathrm{SL}_\Lambda(M)'&\simeq\operatorname{Aut}^{+}(SP_M)\\
\mathrm{GL}_\Lambda(M)^{\prime\deleted{\unsure{1}}\pm}&\simeq\operatorname{Aut}(SP_M)
\end{aligned}\right.\quad,\tag*{(52)}
\]

\noindent\vrule\hspace{.5em}\parbox{.92\linewidth}{et $SP_M$ est un $2$-polyèdre \unsure{triangulé|triangulé-} orienté connexe, régulier du point de vue orienté aussi bien que du point de vue non orienté.}

Calculons le nb de sommets, arêtes, faces de $SP_M$, dans le cas où

\[
\left\{\begin{aligned}
&\Lambda=\mathbb{Z}/n\mathbb{Z}=\prod_{\alpha}\mathbb{Z}/p_\alpha^{\nu_\alpha}\mathbb{Z}\\
&\text{où}\quad n=\prod p_\alpha^{\nu_\alpha}\qquad\text{les }p_\alpha\text{ premiers.}
\end{aligned}\right.\tag*{(53)}
\]

d'où

% archive sheet number: 23
